\VerdierPhysicalPageAnchor{0058}%
\VerdierSemanticAnchor{ch1-1-5-10-converse}
\noindent\textbf{Réciproquement :}\par

\VerdierSemanticAnchor{ch1-1-5-10}
\VerdierSemanticAnchor{ch1-1-5-10-proposition}
\noindent\textbf{Proposition 1.5.10.}\quad
\emph{Tout morphisme $u\colon F_1\to F_2$ entre les foncteurs
sous-jacents, tels que pour tout $i\in[1,n]$ le diagramme (1.5.9.2)
soit $\beta(e_i)$-commutatif, est un $\beta$-morphisme de foncteurs.}

\VerdierSemanticAnchor{ch1-1-5-10-proof}
En effet, pour tout $\sigma\in\mathbb Z^n$, désignons par $R(\sigma)$
la propriété :

\VerdierSemanticAnchor{ch1-1-5-10-property}
\noindent\emph{$(R(\sigma))$\qquad Le diagramme ci-après est
$\beta(\sigma)$-commutatif :}

\VerdierSemanticAnchor{ch1-1-5-10-eq1}
\begin{equation}
  \resizebox{0.78\textwidth}{!}{%
    $\begin{tikzcd}[ampersand replacement=\&,column sep=6em,row sep=4em]
      F_1T(\sigma)
        \arrow[r,"{p_1(\sigma)}"]
        \arrow[d,"{u*T(\sigma)}"']
        \& T'(\sigma)F_1
          \arrow[d,"{T'(\sigma)*u}"] \\
      F_2T(\sigma)
        \arrow[r,"{p_2(\sigma)}"]
        \& T'(\sigma)F_2
    \end{tikzcd}$%
  }\;.
  \tag*{\ensuremath{(1.5.10.1)_{\sigma}}}
\end{equation}

L'examen des différents diagrammes
$(1.5.10.1)_{\sigma_1}$, $(1.5.10.1)_{\sigma_2}$,
$(1.5.10.1)_{\sigma_1+\sigma_2}$ montre immédiatement qu'on a :

\VerdierSemanticAnchor{ch1-1-5-10-closure}
\noindent\emph{Si deux des propriétés $R(\sigma_1)$, $R(\sigma_2)$,
$R(\sigma_1+\sigma_2)$ sont vraies, la troisième l'est aussi.}

\VerdierSemanticAnchor{ch1-1-5-10-conclusion}
De plus par hypothèse, on a les propriétés $R(e_i)$, $i\in[1,n]$ et
$R(0)$. On en déduit immédiatement que la propriété $R(\sigma)$ est
vraie pour tout $\sigma\in\mathbb Z^n$.

\VerdierSemanticAnchor{ch1-1-5-11}
\noindent\textbf{1.5.11.}\quad
Une $\mathbb Z^n$-catégorie stricte est une catégorie munie d'une
représentation de $\mathbb Z^n$ dans son groupe des automorphismes. Pour
définir une telle représentation, il suffit bien entendu de se donner
$n$ automorphismes :
$T(e_1),\ldots,T(e_n)$ qui commutent deux à deux.

\VerdierSemanticAnchor{ch1-1-5-12}
\noindent\textbf{1.5.12.}\quad
Le problème consistant à déterminer sur une catégorie donnée les
structures de $\mathbb Z^n$-catégorie (non nécessairement stricte) est
beaucoup moins simple. Ainsi, si on se borne à déterminer, sur une
catégorie $\mathcal C$, les structures de $\mathbb Z^n$-catégorie
telles que pour tout $\sigma\in\mathbb Z^n$ on ait
$T(\sigma)=\text{identité de }\mathcal C$, on a alors à se donner,
pour tout couple $\sigma_1,\sigma_2$ d'éléments de $\mathbb Z^n$, un
élément $c(\sigma_1,\sigma_2)$ du groupe des automorphismes du foncteur
identique de $\mathcal C$ (ce groupe est commutatif), cette donnée étant
soumise à vérifier les identités (1.2.1.3). Par suite, pour se donner une
telle structure, il faut se donner un 2-cocycle normalisé non homogène de
$\mathbb Z^n$ à valeurs dans le groupe des automorphismes du foncteur
identique de $\mathcal C$. On vérifie d'ailleurs immédiatement que deux
telles structures sont isomorphes, par un $\mathbb Z^n$-foncteur
$(F,m)$ tel que $F=\text{identité de }\mathcal C$, si et seulement si
les cocycles correspondants sont dans la même classe de cohomologie.
